% !TeX root = ../main.tex
\section{Controlled Verification Results}

\begin{table}[t]
  \caption{D1 counts on 20 co-developed patches. The five changed nodes and 12 changed edges are regression items, not independent samples of general accuracy; full-graph counts repeat baseline structure.}
  \label{tab:detection-results}
  \centering
  \small
  \begin{tabular}{lrrr}
    \hline
    Measure & Matched & Extra & Missed \\
    \hline
    Full-graph nodes & 105 & 0 & 0 \\
    Full-graph edges & 108 & 0 & 0 \\
    Changed nodes & 5 & 0 & 0 \\
    Changed edges & 12 & 0 & 0 \\
    \hline
  \end{tabular}
\end{table}

\begin{table}[t]
  \caption{Decision, evidence, and replay outcomes on controlled D1 cases.}
  \label{tab:outcome-results}
  \centering
  \small
  \begin{tabular}{lrr}
    \hline
    Measure & Matches & Total \\
    \hline
    Classification & 20 & 20 \\
    Violation rule set & 7 & 7 \\
    Call-site file and line & 9 & 9 \\
    Missing-relation anchor file and line & 2 & 2 \\
    Deterministic cases & 20 & 20 \\
    \hline
  \end{tabular}
\end{table}

\begin{table}[t]
  \caption{P3 Git-diff gate outcomes and measured scope on the D1 replay.}
  \label{tab:phase3-results}
  \centering
  \small
  \begin{tabular}{lr}
    \hline
    Measure & Result \\
    \hline
    Classification agreement & 20/20 \\
    Merge-decision agreement & 20/20 \\
    Changed-file agreement & 20/20 \\
    Architecture-delta agreement & 20/20 \\
    Incremental/full-scan agreement & 20/20 \\
    Violation rule-set agreement & 7/7 \\
    Call-site exact line & 9/9 \\
    Anchor exact line & 2/2 \\
    Cache hit on repeated check & 20/20 \\
    Normalized cold/warm replay & 20/20 \\
    Incremental files / head files & 57/189 \\
    Cold p50 / p95 (ms) & 518.51 / 531.05 \\
    Warm p50 / p95 (ms) & 242.62 / 249.30 \\
    \hline
  \end{tabular}
\end{table}

\begin{table*}[t]
  \caption{D2 within-tool regression confusion matrices for the frozen v0.1 and provenance-aware v0.2 analyzers on the same co-developed 40-signal corpus.}
  \label{tab:pattern-results}
  \centering
  \small
  \begin{tabular}{lrrrrrrrr}
    \hline
    & \multicolumn{4}{c}{v0.1} & \multicolumn{4}{c}{v0.2} \\
    Detector & TP & FP & FN & TN & TP & FP & FN & TN \\
    \hline
    \texttt{fetch} & 4 & 0 & 0 & 4 & 4 & 0 & 0 & 4 \\
    PostgreSQL & 4 & 1 & 0 & 3 & 4 & 0 & 0 & 4 \\
    Redis & 2 & 0 & 2 & 4 & 4 & 0 & 0 & 4 \\
    AMQP publish & 4 & 2 & 0 & 2 & 4 & 0 & 0 & 4 \\
    AMQP consume & 4 & 1 & 0 & 3 & 4 & 0 & 0 & 4 \\
    \hline
    Overall & 18 & 4 & 2 & 16 & 20 & 0 & 0 & 20 \\
    \hline
  \end{tabular}
\end{table*}

\textbf{RQ1: graph reconstruction and detector signals.}
The unchanged repository was reconstructed as five components and five relationships with a no-impact classification. Table~\ref{tab:detection-results} reports no false positive or false negative in D1. Full-graph totals aggregate 105 expected nodes and 108 expected edges across 20 patched repositories. Only five node changes and 12 edge changes contribute positive changed-graph items. These finite, co-developed cases test contract agreement; they do not estimate performance on unseen projects.

Table~\ref{tab:pattern-results} gives the signal-level regression result. On D2, v0.1 achieved precision 0.818, recall 0.900, F1 0.857, and specificity 0.800. Its four false positives were a PostgreSQL-shaped local query, two AMQP-publish lookalikes, and an AMQP-consume lookalike; its two false negatives were Redis \texttt{hSet} and \texttt{lPush}. On the same annotated signals, v0.2 matched all 20 positives and rejected all 20 annotated negatives (precision, recall, F1, and specificity each 1.000 on this corpus). Package-provenance tracking removed the four method-name false positives, while the expanded Redis operation set recovered the two missed positives. Because D2 was constructed from known v0.1 weaknesses, this is targeted regression evidence and not an estimate over arbitrary TypeScript syntax or an external baseline comparison.

\textbf{RQ2: change classification.}
ArchSync matched all 20 D1 labels: nine no-impact, seven violation, and four evolution cases. It also matched the complete expected rule-identifier set for all seven violations. The added allow and require-path cases exercise rule semantics beyond direct deny and require checks. Redis, AMQP, Fraud Service, and payment-cache additions were classified as evolution rather than violation, preserving a human-review path.

\textbf{RQ3: evidence localization.}
All 11 finding-bearing D1 cases are split post hoc into call-site file/line agreement in 9/9 cases and anchor agreement in 2/2 cases, for missing-edge case 08 and missing-path case 17. Both P2 and P3 satisfy the separate checks. These are case-level outcomes: case 15 has two violated rules but contributes one case. An anchor is not a localized absent call, and neither measure evaluates ranking, all supporting locations, or reviewer effort.

\textbf{RQ4: determinism and reproducibility.}
The D1 baseline and all 20 cases produced byte-identical normalized JSON across duplicate executions, giving 42 current-version executions with no replay mismatch. Both D2 analyzer versions were also deterministic across their duplicate runs. P3 produced identical normalized cold and warm results for 20/20 cases, matched all 20 declared architecture deltas, matched a separate full head scan in all 20 cases, and hit the baseline cache on every repeated check. As Table~\ref{tab:phase3-results} shows, component-incremental analysis parsed 57 of 189 TypeScript file instances, $r_{\mathrm{parse}}=0.3016$. On the recorded Windows environment, warm nearest-rank p50 latency was 242.62 ms compared with 518.51 ms cold, a descriptive 53.2\% reduction; warm p95 was 249.30 ms compared with 531.05 ms cold. Local verification reproduced all stored results and repository gates. GitHub Actions independently reproduced the functional benchmark, detector, and Git-diff evidence on both \texttt{ubuntu-latest} and \texttt{windows-latest}; timing values were not treated as cross-platform acceptance thresholds.

\textbf{Boundary of the reported evidence.}
The observations above come from D1, D2, and the P3 reuse of D1. No external-tool result is accepted into this evaluation, and no common label mapping has been established for comparative scoring. The internal v0.1--v0.2 detector comparison and incremental/full-scan equivalence therefore answer different questions from an external baseline: the former exposes targeted regression behavior, while the latter checks whether scoped execution preserves this implementation's full-scan result.

Exact agreement with these oracles does not establish that the expected architecture is complete or appropriate for another system. In particular, a missed relationship outside the detector vocabulary can be absent from both incremental and full scans without producing an equivalence mismatch. Similarly, duplicate executions can reproduce the same incorrect inference. The graph, evidence-location, decision, and replay checks are complementary consistency checks; they do not substitute for independently authored semantic labels or a separately evaluated tool.
